\documentclass[12pt,a4paper]{extarticle}
\usepackage[utf8]{inputenc}
\usepackage[T5]{fontenc}
\usepackage{geometry}
\geometry{a4paper, left=1.5cm, right=1cm, top=1.5cm, bottom=1.5cm, headheight=20pt, headsep=15pt, footskip=28pt}
\usepackage{enumitem}
\usepackage{amsmath}
\usepackage{amssymb}
\usepackage{diagbox}
\usepackage[most]{tcolorbox}
\newcommand{\khongxacdinh}{\text{KXĐ}}
\newcommand{\blackbox}[1]{\texorpdfstring{\fcolorbox{black}{black}{\textcolor{white}{#1}}}{#1}}
\usepackage{tikz}
\definecolor{bookmagenta}{RGB}{58,90,172}
\definecolor{book}{RGB}{58,90,172}
\definecolor{myred}{RGB}{200,40,40}
\newcommand{\bluecheck}{\textcolor{myblue}{\checkmark}}
\newtcolorbox{formulaBox}{colback=gray!5,colframe=bookmagenta,boxrule=0.8pt,arc=2mm}
\newtcolorbox{notebox}{colback=blue!3,colframe=myblue,boxrule=0.8pt,arc=2mm}
\usetikzlibrary{patterns, patterns.meta} % Thêm patterns.meta vào đây
% Gọi package nằm trong thư mục con template
\usepackage{../template/mngocbox}

\begin{document}

% Sử dụng ngoặc vuông [...] để thay đổi linh hoạt các thông số
\makeheadbox[headerright={Tổng ôn 10-11},chude={CHỦ ĐỀ 10},tieude={Hình học không gian},dinhdang={Hình học},lop={12 },gv={Nguyễn Lê Minh Ngọc}]

\vspace{2em}
\makeheading{III}{Bài tập Trắc nghiệm}

\begin{enumerate}
    % Câu 1
    \item \textbf{Câu 1:} Trong hình học không gian
    \begin{enumerate}
        \item Qua ba điểm xác định một và chỉ một mặt phẳng.
        \item Qua ba điểm phân biệt xác định một và chỉ một mặt phẳng.
        \item Qua ba điểm phân biệt không thẳng hàng xác định một mặt phẳng.
        \item Qua ba điểm phân biệt không thẳng hàng xác định một và chỉ một mặt phẳng.
    \end{enumerate}
    % Câu 2
    \item \textbf{Câu 2:} Trong mặt phẳng $(\alpha)$, cho tứ giác $ABCD$ có $AB$ cắt $CD$ tại $E$, $AC$ cắt $BD$ tại $F$, $S$ là điểm không thuộc $(\alpha)$. Giao tuyến của $(SAB)$ và $(SCD)$ là:
    \begin{multicols}{4}
        \begin{enumerate}
            \item $SF$.
            \item $SD$.
            \item $AC$.
            \item $SE$.
        \end{enumerate}
    \end{multicols}
    % Câu 3
    \item \textbf{Câu 3:} Cho hình chóp $S.ABCD$ có đáy là hình thang $ABCD$ ($AB \parallel CD$). Khẳng định nào sau đây \textbf{sai}?
    \begin{enumerate}
        \item Hình chóp $S.ABCD$ có 4 mặt bên.
        \item Giao tuyến của hai mặt phẳng $(SAC)$ và $(SBD)$ là $SO$ ($O$ là giao điểm của $AC$ và $BD$).
        \item Giao tuyến của hai mặt phẳng $(SAD)$ và $(SBC)$ là $SI$ ($I$ là giao điểm của $AD$ và $BC$).
        \item Giao tuyến của hai mặt phẳng $(SAB)$ và $(SAD)$ là đường trung bình của $ABCD$.
    \end{enumerate}
    % Câu 4
    \item \textbf{Câu 4:} Cho hình chóp $S.ABCD$, đáy $ABCD$ là hình bình hành $ABCD$ tâm $O$. Giao tuyến của hai mặt phẳng $(SAC)$ và $(SAD)$ là
    \begin{multicols}{4}
        \begin{enumerate}
            \item $SO$.
            \item $SD$.
            \item $SA$.
            \item $SB$.
        \end{enumerate}
    \end{multicols}
    % Câu 5
    \item \textbf{Câu 5:} Cho tứ diện $ABCD$, $G$ là trọng tâm tam giác $BCD$. Giao tuyến của $(ACD)$ và $(GAB)$ là
    \begin{enumerate}
        \item $AM$ (với $M$ là trung điểm $AB$).
        \item $AN$ (với $N$ là trung điểm $CD$).
        \item $AK$ (với $K$ là hình chiếu của $C$ trên $BD$).
        \item $AH$ (với $H$ là hình chiếu của $B$ trên $CD$).
    \end{enumerate}
    % Câu 6
    \item \textbf{Câu 6:} Cho bốn điểm $A,B,C,D$ không đồng phẳng. Gọi $I,K$ lần lượt là trung điểm hai đoạn thẳng $AD$ và $BC$. $IK$ là giao tuyến của cặp mặt phẳng nào sau đây?
    \begin{multicols}{2}
        \begin{enumerate}
            \item $(IBC)$ và $(KBD)$.
            \item $(IBC)$ và $(KCD)$.
            \item $(IBC)$ và $(KAD)$.
            \item $(ABI)$ và $(KAD)$.
        \end{enumerate}
    \end{multicols}
    % Câu 7
    \item \textbf{Câu 7:} Cho tứ diện $ABCD$. Lấy điểm $M$ sao cho $AM = 2CM$ và $N$ là trung điểm $AD$. Gọi $O$ là một điểm thuộc miền trong của $\Delta BCD$. Giao điểm của $BC$ với $(OMN)$ là giao điểm của $BC$ với
    \begin{multicols}{2}
        \begin{enumerate}
            \item $OM$.
            \item $MN$.
            \item $A,B$ đều đúng.
            \item $A,B$ đều sai.
        \end{enumerate}
    \end{multicols}
    % Câu 8
    \item \textbf{Câu 8:} Cho hình chóp $S.ABCD$ đáy $ABCD$ là hình bình hành tâm $O$. Gọi $M$ là trung điểm của $SB$. Giao điểm của $DM$ và $(SAC)$ là
    \begin{enumerate}
        \item Giao điểm của $DM$ và $SA$.
        \item Giao điểm của $DM$ và $SC$.
        \item Giao điểm của $DM$ và $SO$.
        \item Giao điểm của $DM$ và $BD$.
    \end{enumerate}
\end{enumerate}



\end{document}
